% Auto-extracted table snippet; requires \usepackage{booktabs}.
% Caption: Table 3: Quantitative results of different loss combinations. Here, ‘1’ indicates the use of this loss function.
\begin{tabular}{lllll}
\toprule
 & JPEG ARCNN DnCNN RNAN FBCNN Zhao et al. Ours GT Figure 4: Visual comparisons on the ‘Classic5’ dataset. No. λ λ λ PSNR↑/ SSIM↑ char dct ssim 1 1 1 0 34.42 / 0.906 2 1 0 1 34.39 / 0.907 3 1 1 1 34.46 / 0.908 Table 3: Quantitative results of different loss combinations. &  &  &  \\
\midrule
No. λ λ char 1 1 2 1 3 1 Table 3: Quantitative re Here, ‘1’ indicates the Twitter ARCNN FBCNN Ours Method ARCNN Dn Param (M) 0.11 0 Figure 5: We perform visual comparisons on real-world PSNR (dB) 29.03 29 images from the ‘Twitter’ dataset, employing residual im- ages for visualization. It can be seen that our method yields Table 4: Comparison o smaller residuals in edge details. Q10 across different m Method LB GB CFSSB SASSB PSNR↑/ SSIM↑ Baseline ✓ 34.21 / 0.905 tative outcomes, presen Model-1 ✓ ✓ ✓ 34.33 / 0.906 loss functions play a ro Model-2 ✓ ✓ 34.29 / 0.906 & No. λ λ char 1 1 2 1 3 1 Table 3: Quantitative re &  &  &  \\
 &  &  &  & use of this loss function. \\
 &  &  & ARCNN Dn & CNN RDN FBCNN Zhao et al. Ours \\
 &  &  & 0.11 0 29.03 29 & .56 22.12 71.9 12.18 16.60 .40 30.03 30.12 30.12 30.25 \\
 &  & ✓ ✓ ✓ ✓ ✓ ✓ &  &  \\
Model-3 Ours &  & ✓ ✓ ✓ &  &  \\
 &  & ✓ ✓ ✓ ✓ &  &  \\
\bottomrule
\end{tabular}
